% !TEX root = ../main.tex
\begin{table}[htbp]
\centering
\small
\caption{Tecnologie principali e versioni dichiarate nella configurazione del progetto.}
\label{tab:tecnologie}
\begin{tabularx}{\textwidth}{@{}p{0.30\textwidth}p{0.15\textwidth}X@{}}
\toprule
\textbf{Tecnologia} & \textbf{Versione} & \textbf{Ruolo nel progetto} \\
\midrule
Java & 17 & Livello di compatibilità del codice compilato. \\
Maven & Non fissata & Build, gestione delle dipendenze ed esecuzione dei plugin. \\
PostgreSQL & 16 & DBMS configurato per l'ambiente Docker Compose. \\
Driver JDBC PostgreSQL & 42.7.7 & Accesso al database dalle implementazioni dei repository. \\
HikariCP & 5.1.0 & Gestione del pool di connessioni al database. \\
Flyway & 9.22.3 & Applicazione delle migrazioni versionate dello schema. \\
JUnit Jupiter & 5.10.2 & Definizione ed esecuzione dei test automatici. \\
Mockito & 5.12.0 & Simulazione delle dipendenze nei test isolati. \\
JaCoCo & 0.8.15 & Raccolta e generazione dei report di copertura del codice. \\
\bottomrule
\end{tabularx}
\par\smallskip
\begin{minipage}{\textwidth}
\footnotesize
Le versioni delle dipendenze e la release Java sono dichiarate in \texttt{pom.xml};
la versione principale di PostgreSQL è specificata in \texttt{docker-compose.yml}.
La release Java non identifica necessariamente la versione della JDK usata per la build.
\end{minipage}
\end{table}
